\section{Finite sums, counting and explicit thresholds}

\begin{lemma}[A finite sum separated by two windows]
    \label{thm:chain_two_window_sum}
    \lean{Finset.sum_range_eq_five_parts_of_two_windows,
        Fin.sum_univ_eq_five_parts_of_two_windows}
    \leanok
    Let $M$ be an abelian monoid written additively, and let
    $a,b,\ell,N$ be natural numbers with $a+2\ell\le b$ and
    $b+2\ell\le N$. For a sequence $f:\mathbb N\to M$,
    \begin{align}
        \sum_{t<N}f(t)
         & =\sum_{t<a}f(t)+\sum_{r<2\ell}f(a+r)\notag   \\
         & \quad+\sum_{r<b-(a+2\ell)}f(a+2\ell+r)\notag \\
         & \quad+\sum_{r<2\ell}f(b+r)
        +\sum_{r<N-(b+2\ell)}f(b+2\ell+r).
        \label{eq:chain_two_window_sum}
    \end{align}
    The same identity holds when $f$ is given only on
    $\{0,\ldots,N-1\}$: every shifted index in~(\ref{eq:chain_two_window_sum})
    belongs to this set. Empty windows and empty complementary intervals
    are permitted. This is the additive grouping used
    in~\cite[Section~4.1]{Arad2016RigorousRG}.
\end{lemma}

\begin{proof}\leanok
    Split the sum successively at $b+2\ell$, $b$, $a+2\ell$ and $a$.
    The inequalities on the endpoints give the complementary lengths
    in~(\ref{eq:chain_two_window_sum}). For a finite array, extend $f$ by zero
    outside its domain and reindex each interval by its finite length.
    Each shifted index is less than $N$, so the extension agrees with $f$
    in every summand.
\end{proof}

\begin{lemma}[Two windows and three intervals of nearest-neighbour terms]
    \label{thm:actual_two_window_alphabet_sum}
    \lean{Fin.sum_univ_eq_actual_two_window_alphabet}
    \leanok
    Let $a+2\ell\le b$ and $b+2\ell\le L-1$, and let $h_i$, for
    $0\le i<L-1$, be values in an additive commutative monoid. For the cuts
    $(k_0,\ldots,k_5)=(0,a+1,a+2\ell,b+1,b+2\ell,L)$ put
    $A_j=\sum_{k_j\le i,\ i+1<k_{j+1}}h_i$, and put $W_r^{\mathrm L}=h_{a+r}$
    and $W_r^{\mathrm R}=h_{b+r}$ for $0\le r<2\ell$. Then
    \begin{align}
        \sum_{i=0}^{L-2}h_i
        =A_0+\sum_{r=0}^{2\ell-1}W_r^{\mathrm L}
        +A_2+\sum_{r=0}^{2\ell-1}W_r^{\mathrm R}+A_4.
        \notag
    \end{align}
    Empty windows are included.
\end{lemma}

\begin{proof}\leanok
    Each finite window enumerates precisely its interval of indices.
    Partition all indices into the left interval, the first complete window,
    the middle interval, the second complete window and the right interval,
    and sum in each part. This retains the original alphabet
    of~\cite[eq:H-mid-cut, lines 810--821]{Arad2016RigorousRG}.
\end{proof}

\begin{lemma}[Frequency vectors with two prescribed coordinates]
    \label{thm:chain_two_fixed_frequencies}
    \lean{Nat.card_fun_sum_eq_and_two_eq,
        Nat.card_fun_sum_eq_and_two_eq_of_lt,
        Nat.card_fun_sum_eq_and_two_eq_of_total}
    \leanok
    Let $I$ be a finite set, let $u,v\in I$ be distinct, and let
    $r,i,j$ be natural numbers. Define
    \begin{align}
        \mathcal F_{r,i,j}
         & =\left\{\alpha:I\to\mathbb N:
        \sum_{a\in I}\alpha(a)=r,\quad
        \alpha(u)=i,\quad\alpha(v)=j\right\}.\label{eq:chain_fixed_frequency_set}
    \end{align}
    If $r<i+j$, then the set in~(\ref{eq:chain_fixed_frequency_set}) is empty.
    If $|I|=s+3$ for a natural number $s$, then
    \begin{align}
        |\mathcal F_{r,i,j}|
         & =\begin{cases}
                \displaystyle\binom{r-i-j+s}{s}, & i+j\le r, \\
                0,                               & r<i+j.
            \end{cases}\label{eq:chain_fixed_frequency_count}
    \end{align}
    Taking $r=m+i+j$ in~(\ref{eq:chain_fixed_frequency_count}) gives
    $|\mathcal F_{m+i+j,i,j}|=\binom{m+s}{s}$
    for every natural number $m$.
    The grouped alphabet in~\cite[Section~4.1]{Arad2016RigorousRG}
    has two windows of $2\ell$ labels and three bulk labels;
    for this alphabet the formula applies with $s=4\ell$.
\end{lemma}

\begin{proof}\leanok
    For every $\alpha:I\to\mathbb N$,
    $\alpha(u)+\alpha(v)\le\sum_{a\in I}\alpha(a)$,
    proving the empty case.
    When the total is $m+i+j$, restricting to $I\setminus\{u,v\}$
    is a bijection between $\mathcal F_{m+i+j,i,j}$ and the
    nonnegative integer functions on $s+1$ labels whose sum is $m$.
    Its inverse reinserts $i$ at $u$ and $j$ at $v$.
    Identifying these functions with multisets of cardinality $m$
    gives $\binom{m+s}{s}$.
    When $i+j\le r$, substituting $m=r-i-j$
    gives~(\ref{eq:chain_fixed_frequency_count}).
\end{proof}

\begin{lemma}[Uniform bound for frequency slices]
    \label{thm:chain_fixed_frequency_uniform_bound}
    \lean{Nat.card_fun_sum_eq_and_two_eq_le}
    \leanok
    \uses{thm:chain_two_fixed_frequencies}
    Let $I$ have cardinality $s+3$, and let $u,v\in I$ be distinct.
    For natural numbers $r,i,j$, let $\mathcal F_{r,i,j}$ be the set
    in~(\ref{eq:chain_fixed_frequency_set}). Then
    \begin{align}
        |\mathcal F_{r,i,j}| & \le (r+1)^s.
        \label{eq:chain_fixed_frequency_uniform_bound}
    \end{align}
\end{lemma}

\begin{proof}\leanok
    \uses{thm:chain_two_fixed_frequencies}
    If $r<i+j$, the set is empty. Otherwise put $m=r-i-j$.
    The exact count in~(\ref{eq:chain_fixed_frequency_count}) and the
    elementary binomial estimate give
    \begin{align}
        |\mathcal F_{r,i,j}|=\binom{m+s}{s}
         & \le (m+1)^s\le(r+1)^s.
        \label{eq:chain_frequency_binomial_bound}
    \end{align}
    Combining~(\ref{eq:chain_frequency_binomial_bound}) with the empty case
    proves~(\ref{eq:chain_fixed_frequency_uniform_bound}).
\end{proof}

\begin{lemma}[Frequency slices through a prescribed degree]
    \label{thm:chain_fixed_frequency_degree_sum}
    \lean{Nat.sum_card_fun_sum_eq_and_two_eq_le}
    \leanok
    \uses{thm:chain_fixed_frequency_uniform_bound}
    Let $I$ be a finite set of cardinality $s+3$, with $s\in\mathbb N$,
    and let $u,v\in I$ be distinct. For natural numbers $k,i,j$, use
    $\mathcal F_{r,i,j}$ from~(\ref{eq:chain_fixed_frequency_set}). Then
    \begin{align}
        \sum_{r=0}^{k}|\mathcal F_{r,i,j}|
         & \le (k+1)^{s+1}.
        \label{eq:chain_frequency_degree_sum}
    \end{align}
\end{lemma}

\begin{proof}\leanok
    \uses{thm:chain_fixed_frequency_uniform_bound}
    For each $0\le r\le k$,
    Lemma~\ref{thm:chain_fixed_frequency_uniform_bound} gives
    $|\mathcal F_{r,i,j}|\le(r+1)^s\le(k+1)^s$.
    Summing over the $k+1$ degrees gives
    $(k+1)(k+1)^s=(k+1)^{s+1}$.
\end{proof}

The following scalar estimates are auxiliaries for the two-window argument
of~\cite[Proposition~5.7]{Arad2016RigorousRG}. For natural
$k,d,\ell,b$, write
\begin{align}
    C_{k,d,\ell,b}
     & =(2\ell)^2(b+1)^2(k+1)^{4\ell+1}\notag            \\
     & \qquad\cdot\bigl(2(k+1)\bigr)^{4b+2}d^{4b+4\ell}.
    \label{eq:chain_two_window_count_factor}
\end{align}

\begin{lemma}[Summing the two-window counting factors]
    \label{thm:chain_two_window_count_sum}
    \lean{Nat.two_window_rank_sum_le_factor}
    \leanok
    Let $k,d,\ell,b\in\mathbb N$ with $d\ge1$ and $b\le k$. Then
    \begin{align}
         & (2\ell)^2
        \sum_{\substack{0\le r\le k                    \\0\le i,j\le b}}
        (r+1)^{4\ell}
        \binom{r+i+1}{2i+1}\binom{r+j+1}{2j+1}
        d^{2(i+j)+4\ell}\notag       \\
         & \qquad\le C_{k,d,\ell,b}.
        \label{eq:chain_two_window_count_sum}
    \end{align}
    The binomial factors retain the upper arguments
    $r+i+1$ and $r+j+1$ printed in~\cite[Proposition~5.7]{Arad2016RigorousRG};
    the count of commuting gaps is recorded
    in~\cite{gap:aklv13_interpolation_and_word_count}.
\end{lemma}

\begin{proof}\leanok
    For $r\le k$ and $i\le b\le k$, the elementary binomial estimate gives
    \begin{align}
        \binom{r+i+1}{2i+1}
         & \le (r+i+1)^{2i+1}
        \le\bigl(2(k+1)\bigr)^{2b+1}.
        \label{eq:chain_two_window_gap_bound}
    \end{align}
    Apply~(\ref{eq:chain_two_window_gap_bound}) to both binomials.
    Also $(r+1)^{4\ell}\le(k+1)^{4\ell}$ and
    $d^{2(i+j)+4\ell}\le d^{4b+4\ell}$, since $d\ge1$.
    There are $(k+1)(b+1)^2$ triples. Multiplying by this number and by
    $(2\ell)^2$ gives the product~(\ref{eq:chain_two_window_count_factor}).
\end{proof}

\begin{lemma}[An explicit bound for the counting product]
    \label{thm:chain_two_window_count_power}
    \lean{Nat.two_window_rank_factor_le_pow}
    \leanok
    Let $k,d,\ell,b\in\mathbb N$ with $k\ge1$, $d\ge2$,
    $\ell\ge1$ and $b\le k$. For the product
    in~(\ref{eq:chain_two_window_count_factor}),
    \begin{align}
        C_{k,d,\ell,b}
         & \le(dk)^{17(\ell+b)}.
        \label{eq:chain_two_window_count_power}
    \end{align}
\end{lemma}

\begin{proof}\leanok
    Put $X=dk\ge2$. The hypotheses give $d\le X$,
    $b+1\le k+1\le X$, $2\ell\le2^\ell\le X^\ell$ and
    $2(k+1)\le X^2$. Substituting these bounds
    into~(\ref{eq:chain_two_window_count_factor}) yields
    $C_{k,d,\ell,b}\le X^{10\ell+12b+7}\le X^{17(\ell+b)}$, since $\ell\ge1$.
\end{proof}

\begin{definition}[Ordinary labels]
    \label{def:two_marker_ordinary_labels}
    \lean{MvPolynomial.twoMarkerOrdinaryLabels}
    \leanok
    For a finite alphabet $I$ and distinguished labels $u,v\in I$, let
    $\mathcal O_{u,v}=I\setminus\{u,v\}$.
    The definition does not require the two labels to be distinct.
\end{definition}

\begin{lemma}[One sample family for every degree and frequency slice]
    \label{thm:two_window_common_scalar_recovery}
    \lean{Polynomial.exists_twoWindow_wordCoefficient_recovery}
    \leanok
    \uses{def:two_marker_ordinary_labels,def:full_word_coefficient}
    Let $\ell,K$ be nonnegative integers, let
    $I=(\{0,\ldots,2\ell-1\}\sqcup\{0,\ldots,2\ell-1\})\sqcup\{0,1,2\}$,
    and choose $0\le p,s<2\ell$.
    Let $u$ be the label $p$ in the first window and $v$ the label $s$ in the
    second window. Let $R$ be a ring and a complex algebra, without any
    commutativity assumption on $R$, and put $N=(K+1)^{4\ell+1}$.
    For $H:I\to R$, $r\ge0$ and $\nu:I\to\N$, let $\mathcal W_{H,r}(\nu)$ be
    the sum of the ordered products $H_{w(0)}\cdots H_{w(r-1)}$ over the words
    $w$ of length $r$ in which each letter $a$ occurs $\nu(a)$ times.

    There are complex sample points $z_t:I\to\C$, $0\le t<N$, with the
    following property. For every $i,j\ge0$ and every scalar mask
    $\mu:(I\to\N)\to\C$, there are complex coefficients $\gamma_t$ such
    that, for every family $H:I\to R$ and every $0\le r\le K$,
    \begin{align}
         & \sum_{\substack{\nu:I\to\N                                   \\ \sum_a\nu(a)=r\\
                   \nu(u)=i,\ \nu(v)=j}}
        \mu(\nu)\,\mathcal W_{H,r}(\nu)\notag                                            \\
         & \qquad=\sum_{t=0}^{N-1}\gamma_t
                                  [x^iy^j]\left(D_t+yH_v+xH_u\right)^r,
        \qquad D_t=\sum_{a\in\mathcal O_{u,v}}z_t(a)H_a.
        \label{eq:two_window_common_scalar_recovery}
    \end{align}
    Here $x,y$ are central formal variables. The selected terms $H_u,H_v$
    remain unweighted. The points precede $i,j,\mu$, while the recovery
    coefficients precede $H$ and every degree $r\le K$.
\end{lemma}

\begin{proof}\leanok
    \uses{def:two_marker_ordinary_labels,def:full_word_coefficient}
    The distinguished labels are distinct, so $\mathcal O=\mathcal O_{u,v}$
    has $4\ell+1$ members. Put $G=\{0,\ldots,K\}^{\mathcal O}$.
    The nodes $0,\ldots,K$ are distinct, so the Vandermonde matrix
    $V_{s,t}=s^t$ is invertible. For $\alpha,g\in G$, define
    $E_{g,\alpha}=\prod_{a\in\mathcal O}g(a)^{\alpha(a)}$ and
    $R_{\alpha,g}=\prod_{a\in\mathcal O}(V^{-1})_{\alpha(a),g(a)}$; then
    $\sum_{g\in G}R_{\alpha,g}E_{g,\beta}=\delta_{\alpha,\beta}$.
    Take $z_g:I\to\C$ to agree with $g$ on $\mathcal O$ and to vanish
    elsewhere.

    Fix $i,j,\mu$. For $\alpha\in G$, let $\nu_{i,j,\alpha}$ be the
    full count vector obtained by assigning $i$ at $u$, $j$ at $v$, and
    $\alpha(a)$ at each ordinary label $a$, and put
    $\gamma_g=\sum_{\alpha\in G}\mu(\nu_{i,j,\alpha})R_{\alpha,g}$.
    For $r\le K$ every ordinary coordinate of a count vector in the slice is
    at most $K$, so restriction to the ordinary labels identifies the slice
    with a subset of $G$, and every slice sum is a zero-extended sum over $G$.
    Substituting the recovery identity and exchanging the finite sums gives
    recovery with the weights $\gamma_g$. The counted-word identity identifies
    each evaluated ordinary monomial sum with $[x^iy^j](D_g+yH_v+xH_u)^r$,
    preserving the order of the factors. The box has exactly $N$ elements;
    reindexing it by $0\le t<N$ gives~(\ref{eq:two_window_common_scalar_recovery}).
\end{proof}

\begin{definition}[Polynomially bounded functions]
    \label{def:nat_polynomial_bounded}
    \lean{IsNatPolynomialBounded,
        IsNatPolynomialBounded.const,
        IsNatPolynomialBounded.id,
        IsNatPolynomialBounded.of_le,
        IsNatPolynomialBounded.add,
        IsNatPolynomialBounded.mul,
        IsNatPolynomialBounded.pow,
        IsNatPolynomialBounded.max_one,
        IsNatPolynomialBounded.exists_one_le}
    \leanok
    A function $f$ from the natural numbers to themselves is polynomially
    bounded when $f(n)\le P\,(n+1)^d$ for all $n\ge1$, for natural numbers $P$
    and $d$ chosen before $n$. Constants and the identity are polynomially
    bounded, and the class is closed under pointwise domination, sums,
    products, powers and taking the maximum with one; the coefficient $P$ can be
    taken positive.
\end{definition}

\begin{lemma}[Sequences with prescribed successive relations]
    \label{lem:dependent_sequence_choice}
    \lean{Nat.exists_dependent_sequence}
    \leanok
    Let $(X_i)_{i\ge0}$ be a family of sets, and let
    $R_i\subseteq X_i\times X_{i+1}$ be a relation for each $i$.
    Suppose every $x\in X_i$ has a successor $y\in X_{i+1}$ with
    $(x,y)\in R_i$. For every prescribed $x_0\in X_0$, there is a sequence
    $(x_i)_{i\ge0}$ beginning at $x_0$ such that
    $(x_i,x_{i+1})\in R_i$ at every index.
\end{lemma}

\begin{proof}\leanok
    Choose for each $i$ a map $g_i:X_i\to X_{i+1}$ satisfying
    $(x,g_i(x))\in R_i$ for every $x\in X_i$.
    Recursively define $x_{i+1}=g_i(x_i)$ from the prescribed initial value.
\end{proof}

\begin{lemma}[Restricting finite sums to nonzero weights]
    \label{lem:finset_sum_subtype_ne_zero}
    \lean{Finset.sum_eq_sum_subtype_ne_zero}
    \leanok
    Let $I$ be finite, let $p_i$ be weights in a type with zero, and let $f_i$
    take values in a commutative additive monoid. If $f_i=0$ whenever $p_i=0$,
    then
    \begin{align}
        \sum_{i\in I}f_i & =\sum_{\substack{i\in I \\p_i\ne0}}f_i.
        \notag
    \end{align}
\end{lemma}

\begin{proof}\leanok
    Split the full sum according to whether $p_i$ vanishes. Every term in the
    zero-weight part vanishes by hypothesis.
\end{proof}

\begin{lemma}[Termwise products under Cauchy--Schwarz]
    \label{lem:norm_sum_le_mul_sqrt_mul_sqrt}
    \lean{Finset.norm_sum_le_mul_sqrt_mul_sqrt}
    \leanok
    If elements $z_i$ of a normed group satisfy
    $\lVert z_i\rVert\le c\,a_ib_i$ for real $a_i,b_i$ and $c\ge0$, then
    \begin{align}
        \Bigl\lVert\sum_iz_i\Bigr\rVert
        \le c\Bigl(\sum_ia_i^2\Bigr)^{1/2}\Bigl(\sum_ib_i^2\Bigr)^{1/2}.
        \notag
    \end{align}
\end{lemma}

\begin{proof}\leanok
    Apply the triangle inequality and then the Cauchy--Schwarz inequality to
    $\sum_ia_ib_i$.
\end{proof}

\begin{lemma}[Expansion in selected multilinear arguments]
    \label{thm:peps_selected_multilinear_expansion}
    \lean{MultilinearMap.map_piecewise_sum_smul}
    \leanok
    Let $T$ be a complex multilinear map, let $S$ be a finite set of argument
    positions, and fix $x_i$ for $i\notin S$. Replacing each $x_i$
    ($i\in S$) by $\sum_{j\in J_i} c_{ij}v_{ij}$ gives
    \begin{align}
        T\bigl(x[S\leftarrow\textstyle\sum_j c_{ij}v_{ij}]\bigr)
         & =\sum_{a\in\prod_{i\in S}J_i}
        \Bigl(\prod_{i\in S}c_{i,a_i}\Bigr)
        T\bigl(x[S\leftarrow v_{i,a_i}]\bigr).
        \notag
    \end{align}
\end{lemma}

\begin{proof}\leanok
    Restricting $T$ to the arguments in $S$ with $x_{S^c}$ fixed gives a
    multilinear map on $\prod_{i\in S}V_i$, and expanding the finite sums by
    multilinearity pulls out $\prod_{i\in S}c_{i,a_i}$.
\end{proof}

\begin{lemma}[Explicit thresholds for elementary limits]
    \label{lem:explicit_thresholds}
    \lean{Real.powerRatioThreshold,
        Real.mul_rpow_le_mul_rpow_of_powerRatioThreshold_le,
        Real.le_rpow_of_powerRatioThreshold_le,
        Real.add_mul_rpow_le_mul_rpow_of_le,
        Real.logPowThreshold,
        Real.mul_log_pow_le_mul_rpow_of_logPowThreshold_le,
        Real.logRpowThreshold,
        Real.mul_rpow_mul_log_rpow_le_mul_rpow_of_logRpowThreshold_le,
        Real.stretchedExpThreshold,
        Real.mul_rpow_le_exp_of_stretchedExpThreshold_le,
        Real.mul_exp_neg_le_rpow_neg_of_stretchedExpThreshold_le,
        Real.one_le_log_natCast_of_three_le}
    \leanok
    Let $C\ge0$ and $\eta>0$. Each of the following inequalities holds for
    every real $x$ beyond an explicit closed-form threshold:
    \begin{enumerate}
        \item $Cx^a\le\eta x^b$ for $a<b$, beyond
              $\max\{1,(C/\eta)^{1/(b-a)}\}$;
        \item $C(\log x)^k\le\eta x^a$ for $a>0$ and natural $k$;
        \item $Cx^u(\log x)^L\le\eta x^v$ for $u<v$ and real $L$;
        \item $Cx^s\le\exp(cx^\theta)$, equivalently
              $C\exp(-cx^\theta)\le x^{-s}$, for $c,\theta>0$.
    \end{enumerate}
    Each statement also holds for natural $n$ beyond the ceiling of the
    threshold, and $\log n\ge1$ for natural $n\ge3$.
\end{lemma}

\begin{proof}\leanok
    The first threshold is obtained by raising $C/\eta\le x^{b-a}$ to the
    power $1/(b-a)$. The second uses $\log x\le x^\varepsilon/\varepsilon$ with
    $\varepsilon=a/(2(k+1))$ and the first threshold at exponent $a/2$. The
    third reduces to the second, rounding $L$ up. The fourth compares $x^s$ with
    the Taylor term $(cx^\theta)^m/m!$ of the exponential, $m=\lfloor s/\theta\rfloor+1$.
\end{proof}

\begin{lemma}[Elementary monotonicity and ceiling bounds]
    \label{thm:homogeneous_ceiling_bound}
    \lean{Nat.ceil_le_add_one_mul_of_le_mul, Real.exp_mul_sub_le_exp_mul_sub}
    \leanok
    For $x,A,s\in\mathbb R$ with $0\le x\le As$ and $s\ge1$,
    \begin{align}
        \lceil x\rceil\le(A+1)s.
        \label{eq:homogeneous_ceiling_bound}
    \end{align}
    For $v\ge0$ and $x\le T$, $\exp(vx-c)\le\exp(vT-c)$.
\end{lemma}

\begin{proof}\leanok
    $\lceil x\rceil\le x+1\le As+1\le As+s=(A+1)s$. The exponential is
    increasing and $vx\le vT$.
\end{proof}
